\begin{afterword}
\summaryhead

The post opens with its thesis: most things in the universe do not have minds. Earlier cultures believed that trees, rocks, rivers and hills had spirits, and the author argues that this was not a stupid belief. People reasoned by analogy from their own bodies: flesh thinks, so why not wood; a flood that hurts the tribe looks like an angry river, as a punch looks like an angry man. Nothing a hunter-gatherer could see ruled this out.

Aristotle thought the brain cooled the blood, Egyptian embalmers threw it away, and Alcmaeon's view that it was the seat of intelligence was only a guess. The author cannot say when the brain's role stopped being a guess, or when anyone first took its tangle of neurons for information processing. The author places the beginning of anthropomorphism's descent into being obviously wrong at three discoveries: that tangle, Darwin's theory of natural selection, and the idea of cognition as computation. With them one can give a reason why a tree, a river, an atom, a puppy or a rock lacks a mind or moral reasoning. The post ends with the Zhuangzi story of the happy fish, and the line ``Now we know.''

\responsehead

The post's claim is modest and mostly sound. It argues that animism was a reasonable inference from what people could see, and that knowledge about brains, evolution and computation later gave reasons against it. It keeps its hedges throughout: the author does ``not know enough history,'' ``haven't had much luck,'' and offers the three discoveries as ``where I would put'' the change.

The account of animism has a long history. Hume gave it in 1757: there is ``an universal tendency among mankind to conceive all beings like themselves,'' so that people ``ascribe malice or good-will to every thing, that hurts or pleases us.'' Hume also named the cure the post reaches, knowledge of ``the minute parts'' of bodies. The post does not mention this history.

Where the post guesses at history, the guesses have old answers. The experiment it proposes, cutting nerves to see whether movement and sensation stop, is close to what Galen did in the second century, when he cut nerves and the spinal cord of animals to see the effect; Herophilus had tied the nerves to motion and sensation earlier still. This supports the post's next step. Galen explained what moved through the nerves as a ``psychic pneuma,'' which is the ``mysterious spirit'' the post says such experiments leave room for.

The post's rock does not fit Aristotle, and Aristotle does not fit the post's history. The rock that wants ``to roll downhill, as Aristotle once thought'' is not Aristotle's: the \textsc{Physics} says that falling bodies do not move themselves, since self-motion ``is a characteristic of life.'' And without microscopes Aristotle held that plants have no power of thinking, on grounds of what plants do, the kind of evidence the post's own tree test counts. So the claim that it took ``really good microscopes'' to see that wood does not think is unsupported as stated. What the later discoveries added was better reasons, which is how the post itself frames its tests.

The tests themselves rest on assumptions the post does not argue: that a mind needs complex information processing, and, for the puppy, what ``current theories of evolutionary psychology'' hold, with no source.

In short: a careful, hedged account of why animism was reasonable and what made it look wrong, which is close to Hume's, with an Aristotle who in fact denied thought to plants and self-motion to rocks.
\end{afterword}
